% !TeX root = ta.tex
% ============================================================================================
% BAB III ANALISIS MASALAH
% Pembagian subbab tidak rigid dan dapat bervariasi. Bab ini minimal berisi analisis kebutuhan
% fungsional dan nonfungsional, analisis berbagai alternatif solusi yang dapat ditawarkan, dan
% metode pemilihan solusi yang diusulkan.
% ============================================================================================
\cleardoublepage
\chapter{ANALISIS MASALAH}
\label{chap:analisis-masalah}

\section{Analisis Kondisi Saat Ini}

Penelitian oleh \textcite{zienius2019direct} menganalisis jalur diagnosis tumor otak melalui pemindaian \textit{Direct Access Computed Tomography} (DACT) di layanan primer. Studi berbasis populasi ini dilakukan selama lima tahun di wilayah Lothian, Skotlandia, dengan mengevaluasi 2.938 rujukan dari dokter umum yang mencurigai adanya patologi intrakranial. Selain menilai keandalan CT kepala sebagai modalitas diagnostik, studi ini mengevaluasi efektivitas dua pedoman rujukan utama, yaitu pedoman Kernick dan NICE 2005, dalam mengidentifikasi pasien yang berisiko mengidap tumor otak.

\subsection{Alur Diagnosis Tumor Otak}
Saat ini, diagnosis tumor otak di layanan primer dilakukan melalui serangkaian tahapan berikut \cite{zienius2019direct}:

\begin{enumerate}
	\item Pasien Merasakan Gejala dan Memeriksakan Diri ke Dokter Umum\\
	Gejala tumor otak bersifat tidak khas dan mudah dikira penyakit lain, sehingga pasien rata rata harus mengunjungi dokter umum sebanyak tiga kali atau lebih sebelum akhirnya mendapat rujukan \cite{zienius2019direct}.
	
	\item Dokter Umum Membuat Rujukan \textit{Direct Access CT} (DACT)\\
	Dokter umum dapat merujuk pasien langsung untuk pemindaian CT kepala tanpa melalui spesialis terlebih dahulu. Namun, pedoman rujukan yang tersedia yaitu pedoman Kernick dan NICE 2005 terbukti tidak akurat dalam memilah pasien, dengan nilai prediksi positif hanya sebesar 3\% dan tingkat kesepakatan antar pedoman yang rendah (\textit{weighted kappa} = 0,33) \cite{zienius2019direct}.
	
	\item Pasien Menunggu Jadwal Pemindaian\\
	Setelah rujukan dibuat, pasien tidak langsung menjalani pemindaian. Rata rata waktu tunggu dalam studi ini adalah 17,3 hari untuk seluruh pasien, dan 14,6 hari untuk pasien yang akhirnya terdiagnosis tumor signifikan \cite{zienius2019direct}. Karena DACT bukan jalur darurat, keterlambatan ini dianggap wajar secara prosedural, namun dalam konteks tumor otak, setiap hari penundaan tetap memiliki implikasi klinis yang nyata.
	
	\item Pemindaian Dilakukan dan Dianalisis oleh Konsultan Neuroradiologi\\
	Seluruh pemindaian CT kepala dilakukan di satu pusat neurosains terpusat dan dilaporkan oleh konsultan neuroradiologi. CT kepala terbukti cukup andal sebagai alat skrining karena tidak ditemukan satu pun hasil negatif palsu dalam studi ini. Namun, \textit{diagnostic yield} untuk tumor signifikan hanya sebesar 1,43\% \cite{zienius2019direct}, yang berarti sebagian besar pemindaian tidak menemukan tumor. Pelaksanaan \textit{neuroimaging} sejak awal turut mempersingkat waktu yang dibutuhkan hingga hasil analisis tersedia bagi klinisi. Hal ini tercermin dalam temuan \textcite{carey2019impact}, yang melaporkan bahwa median waktu menuju diagnosis tumor otak ganas pada kelompok yang menjalani \textit{neuroimaging} awal hanya sebesar 8 hari sejak kunjungan pertama terkait keluhan sakit kepala ($p < 0{,}001$).
	
	\item Laporan Dikembalikan ke Dokter Umum untuk Menentukan Tindak Lanjut\\
	Dari 559 pasien dengan temuan abnormal nontumor, sebanyak 34,2\% tetap perlu dirujuk untuk pencitraan lanjutan atau penilaian spesialis \cite{zienius2019direct}. Ketiadaan standar pelaporan yang baku mempersulit interpretasi dokter umum, khususnya terhadap temuan insidental.
	
	\item Pasien Dirujuk ke Spesialis dan Dikonfirmasi melalui Histopatologi\\
	Apabila hasil pemindaian mencurigakan, pasien dirujuk ke spesialis untuk konfirmasi melalui histopatologi. Dari 42 tumor signifikan yang ditemukan, jenis terbanyak adalah metastasis (40\%), diikuti tumor hipofisis (19\%), meningioma (17\%), \textit{glioblastoma multiforme} (12\%), glioma non GBM (10\%), dan limfoma SSP (2\%) \cite{zienius2019direct}.
\end{enumerate}

Secara keseluruhan, alur DACT berfungsi dengan baik secara teknis dan mampu mendukung diagnosis tumor otak yang lebih cepat melalui \textit{neuroimaging} awal, sebagaimana ditunjukkan oleh \textcite{carey2019impact} dengan median waktu diagnosis sebesar 8 hari sejak kunjungan pertama. Namun, \textcite{zienius2019direct} menegaskan bahwa sistem ini masih lemah pada tahap seleksi pasien dan interpretasi hasil, sehingga efektivitas dan efisiensi alur secara keseluruhan belum optimal dan masih membutuhkan peningkatan lebih lanjut.

\subsection{Penggunaan AI dalam CT Scan Otak}

Meskipun teknologi AI telah dikembangkan untuk membantu interpretasi \textit{non contrast CT head} (NCCTH), penggunaannya saat ini masih berada pada tahap evaluasi berbasis data nyata dan belum memasuki tahap implementasi klinis. Studi AI REACT \cite{fu2024ai} merupakan contoh konkret dari tahap ini: sebuah studi \textit{multireader multicase} yang menguji alat AI bernama \textit{qER EU 2.0} pada 150 pemindaian nyata dari pasien gawat darurat, melibatkan 30 pembaca dari empat rumah sakit NHS di Inggris.

Hasil evaluasi menunjukkan bahwa bantuan AI meningkatkan sensitivitas pembaca secara bermakna dalam mendeteksi abnormalitas kritis, namun disertai penurunan spesifisitas yang signifikan \cite{fu2024ai}. Lebih lanjut, performa algoritma terbukti tidak merata, yakni deteksi \textit{mass effect}, yang merupakan salah satu tanda adanya lesi desak ruang seperti tumor, hanya menghasilkan AUC sebesar 0,604, jauh di bawah subkelompok patologi lainnya \cite{fu2024ai}. Para penulis secara eksplisit menyatakan bahwa hasil studi retrospektif ini tidak dapat langsung diterapkan dalam praktik klinis, dan menegaskan perlunya studi prospektif lanjutan sebelum teknologi ini dapat diimplementasikan secara bertanggung jawab \cite{fu2024ai}.

Penelitian \textcite{khan2024comparative} juga menunjukkan bahwa klasifikasi biner (ada/tidak tumor) memiliki akurasi lebih tinggi dibandingkan multikelas (jenis tumor). Pada CNN kustom misalnya, akurasi biner mencapai 99,62\%, sedangkan multikelas hanya 98,09\%. Pola yang sama juga terlihat pada ResNet18 dan VGG16 \cite{khan2024comparative}. Meski begitu, hal ini wajar karena semakin sedikit kelas yang harus dibedakan, semakin kecil pula peluang model salah menebak, dengan selisih akurasi yang juga tidak terlalu besar (sekitar 0,3\%-1,5\%). Selain itu, klasifikasi multikelas tetap lebih bermanfaat secara klinis karena memberikan informasi jenis tumor, bukan hanya ada atau tidaknya tumor.

\section{Analisis Kebutuhan}
\subsection{Identifikasi Masalah Pengguna}
Terdapat dua kelompok pengguna utama yang menghadapi masalah terkait penerapan AI dalam medis, yaitu profesional medis (dokter dan radiolog) serta pasien.
Bagi profesional medis, tantangan utamanya tidak hanya terletak pada adopsi teknologi baru, tetapi juga pada keterbatasan manusiawi dan beban kerja yang ekstrem dalam praktik klinis. Berdasarkan \textcite{zhang2023diagnostic}, tingkat kesalahan diagnostik dalam radiologi berkisar antara 10-26\% dengan 60-80\% kesalahan  diklasifikasikan sebagai kesalahan perseptual atau kegagalan dalam mendeteksi adanya lesi. Faktor kelelahan turut berperan signifikan dengan survei mencatat 47\% radiolog mengalami \textit{burnout} dan 58\% mengalami gejala fisik seperti kelelahan mata yang dapat mengganggu akurasi interpretasi.

Di Indonesia, masalah ini menjadi jauh lebih kritis akibat ketimpangan rasio tenaga medis. Data analisis pasar layanan radiologi mencatat bahwa Indonesia hanya memiliki sekitar 2.161 radiolog terdaftar untuk melayani populasi sebesar 277 juta jiwa, atau setara dengan rasio 1,2 radiolog per 100.000 penduduk \autocite{insights102023indonesia}. Kondisi kekurangan SDM ini secara langsung meningkatkan beban kerja radiolog, yang berpotensi memperparah risiko \textit{burnout} dan kesalahan diagnosis. Selain itu,
profesional medis juga menghadapi masalah rendahnya kepercayaan terhadap keandalan sistem AI. Masalah ini berakar dari kurangnya edukasi, pelatihan, dan minimnya keterlibatan mereka dalam proses pengembangan sistem \autocite{syaftahan2025transformasi}. Hal ini menciptakan paradoks bahwa teknologi yang sangat dibutuhkan untuk mengurangi \textit{human error} justru diragukan validitasnya oleh penggunanya sendiri.

Dari sisi pasien, tantangan utama yang menghambat adopsi \textit{Artificial Intelligence} (AI) dalam praktik klinis adalah kurangnya kepercayaan terhadap validitas diagnosis yang melibatkan algoritma. Berbeda dengan asumsi umum bahwa teknologi canggih akan meningkatkan rasa aman pasien, literatur terbaru justru menunjukkan fenomena sebaliknya, yaitu keterlibatan AI sering kali dipersepsikan sebagai faktor yang mendegradasi kualitas layanan medis.

Riset eksperimental terbaru yang dilakukan oleh \textcite{chen2025impact} memberikan bukti bahwa kepercayaan pasien terhadap dokter dan intensi untuk melakukan pengobatan mencapai titik tertinggi justru ketika dokter secara eksplisit menyatakan tidak menggunakan AI dalam proses diagnosisnya dengan skor rata-rata kepercayaan 0,63. Sebaliknya, ketika pasien mengetahui bahwa dokter menggunakan AI secara ekstensif, tingkat kepercayaan mereka menurun ke angka 0,30 \autocite{chen2025impact}. Hal ini mengindikasikan bahwa di mata pasien, penggunaan AI belum dianggap sebagai peningkatan kapabilitas diagnostik, melainkan sebagai substitusi yang mengurangi nilai keahlian manusia.

Penurunan kepercayaan ini bukan tanpa alasan. Skeptisisme pasien berakar pada beberapa kekhawatiran fundamental. Pertama, risiko halusinasi AI, seperti model generatif dapat memproduksi informasi medis yang terdengar masuk akal dan meyakinkan tetapi secara faktual salah atau dibuat-buat. Ketakutan akan diagnosis yang keliru akibat kesalahan algoritma ini menjadi penghalang psikologis utama. Kedua, masalah transparansi dan akuntabilitas. Pasien sering kali memandang AI sebagai \textit{black-box} yang proses pengambilan keputusannya tidak dapat dijelaskan secara intuitif, memicu keraguan apakah keputusan medis tersebut didasarkan pada kondisi unik pasien atau sekadar pola statistik semata \autocite{chen2025impact}.

Selain itu, terdapat korelasi negatif antara persepsi profesionalisme dan penggunaan AI. Pasien cenderung mempertanyakan kompetensi dan penilaian klinis dokter yang terlalu bergantung pada alat bantu AI, terutama jika alat tersebut bersifat umum dan tidak dirancang khusus untuk medis. Penggunaan alat semacam ini dianggap mengaburkan batas tanggung jawab profesional dan menurunkan standar interaksi manusiawi yang menjadi inti dari layanan kesehatan \autocite{chen2025impact}. Temuan ini menegaskan bahwa diperlukan transparansi, validasi klinis yang ketat, dan jaminan pengawasan manusia untuk mengatasi ketidakpercayaan pasien.

Berdasarkan analisis dari studi literatur yang dijelaskan pada Bab \ref{sec:penelitian terkait tumor}, analisis kondisi saat ini, dan analisis permasalahan pengguna, diidentifikasikan 4 permasalahan utama, yaitu:

\begin{enumerate}
	\item Seleksi Pasien Kurang Akurat \\
	Pedoman rujukan yang digunakan dokter umum, yaitu pedoman Kernick dan NICE 2005, terbukti memiliki nilai prediksi positif yang sangat rendah, hanya sebesar 3\%, dengan tingkat kesepakatan antar pedoman yang buruk (\textit{weighted kappa} = 0,33). Artinya, kedua pedoman ini hampir tidak lebih baik dari tebakan acak dalam memilah pasien yang benar-benar berisiko tumor. Implikasinya, sebagian besar pemindaian yang dilakukan tidak menemukan tumor (\textit{diagnostic yield} hanya 1,43\%), sehingga sumber daya radiologi terbuang untuk kasus yang tidak relevan, sementara pasien berisiko tinggi bisa saja terlambat tersaring \cite{zienius2019direct}.
	
	\item Keterlambatan Diagnosis \\
	Rata-rata pasien harus mengunjungi dokter umum tiga kali atau lebih sebelum mendapat rujukan, ditambah waktu tunggu pemindaian rata-rata sebesar 17,3 hari \cite{zienius2019direct}. Meskipun prosedur ini dianggap wajar secara administratif, dalam konteks tumor otak setiap hari penundaan berdampak signifikan terhadap prognosis. Studi \textcite{carey2019impact} membuktikan bahwa \textit{neuroimaging} awal mampu mempersingkat waktu diagnosis hingga median 8 hari sejak kunjungan pertama, namun sayangnya akses ke jalur cepat ini masih bergantung pada keputusan dokter umum yang disaring dengan pedoman yang tidak akurat.
	
	\item Ketimpangan Sumber Daya Radiologi \\
	Rasio 1,2 radiolog per 100.000 penduduk di Indonesia merupakan angka yang jauh dari memadai untuk menopang sistem diagnosis berbasis pencitraan \autocite{insights102023indonesia}. Beban kerja yang berlebihan ini tidak hanya menciptakan risiko \textit{burnout} (yang dialami oleh 47\% radiolog secara global), tetapi juga mendegradasi akurasi interpretasi secara sistemik, mengingat 60 hingga 80\% kesalahan diagnostik dalam radiologi bersifat perseptual, yakni radiolog gagal mendeteksi lesi yang sebenarnya ada \cite{zhang2023diagnostic}. Kondisi ini menjadikan keterlibatan AI bukan sekadar pilihan, melainkan kebutuhan yang mendesak.
	
	\item Kurangnya Kepercayaan dari Klinisi dan Pasien \\
	Isu ini berjalan dua arah sekaligus. Di satu sisi, profesional medis meragukan AI karena kurangnya edukasi dan minimnya keterlibatan mereka dalam proses pengembangannya \autocite{syaftahan2025transformasi}. Di sisi lain, pasien justru memberikan kepercayaan tertinggi kepada dokter yang secara eksplisit \textit{tidak} menggunakan AI (skor rata-rata 0,63), dibandingkan dokter yang menggunakannya secara ekstensif (skor 0,30) \autocite{chen2025impact}. Kedua fenomena ini berakar dari sifat \textit{black-box} sistem AI medis, baik klinisi maupun pasien tidak dapat menelusuri dasar suatu keputusan sehingga kepercayaan sulit terbentuk meskipun performa teknisnya tinggi. Karena itu, interpretabilitas perlu menjadi prasyarat pengembangan AI medis, bukan sekadar fitur tambahan. Sistem yang mampu menjelaskan rasionalitas keputusannya secara transparan akan lebih mudah diaudit klinisi sekaligus lebih dipercaya pasien. Tanpa interpretabilitas ini, krisis kepercayaan akan terus berulang terlepas dari secanggih apa pun sistem yang ditawarkan.
\end{enumerate}

Dengan mempertimbangkan keempat permasalahan tersebut, dapat disimpulkan bahwa meskipun alur diagnosis tumor otak melalui jalur \textit{Direct Access CT} telah berfungsi secara teknis, sistem ini masih memiliki berbagai kelemahan yang signifikan dari sisi klinis maupun struktural. Seleksi pasien yang tidak akurat, keterlambatan diagnosis yang berdampak klinis nyata, ketimpangan sumber daya radiologi yang ekstrem, khususnya di Indonesia, serta defisit kepercayaan terhadap AI dari sisi klinisi maupun pasien. Oleh karena itu, tugas akhir ini bertujuan untuk mengembangkan model AI yang dilengkapi dengan modul \textit{interpretability} berupa visualisasi dalam mendeteksi tumor otak melalui citra CT scan kepala.

Model ini diharapkan dapat membuka peluang dalam pembuatan sistem yang memperketat seleksi pasien sejak tahap awal, mempersingkat waktu menuju diagnosis, serta membantu radiolog dalam menganalisis citra secara lebih efisien untuk mengatasi keterbatasan sumber daya manusia. Dengan menyediakan visualisasi penjelas terhadap setiap hasil prediksi, model ini diharapkan dapat menjembatani kesenjangan antara akurasi teknis dan kepercayaan klinis, sekaligus memperluas peluang integrasi dalam pengembangan sistem pendukung keputusan medis yang lebih terbuka, transparan, dan bertanggung jawab. Pengembangan model ini diharapkan memiliki peluang sebagai berikut.

\begin{enumerate}
	\item Pengembangan model AI dengan modul \textit{Explainable AI} berupa visualisasi \textit{heatmap} untuk mengidentifikasi area mencurigakan pada \textit{CT scan}, menggantikan pedoman seleksi pasien yang tidak akurat
	\item Standardisasi hasil pemindaian melalui model AI yang konsisten antar waktu pemeriksaan, mengurangi kesalahan perseptual radiolog akibat kelelahan/beban kerja berlebihan
	\item Sistem pendukung keputusan awal yang mempersingkat waktu dari kunjungan pertama hingga diagnosis, mengurangi keterlambatan yang berdampak klinis pada pasien tumor otak
	\item Implementasi \textit{Explainable AI} sebagai fitur interpretabilitas yang menjelaskan dasar prediksi kepada klinisi dan pasien, menjembatani kesenjangan kepercayaan terhadap AI
\end{enumerate}

\subsection{Kebutuhan Fungsional}
Untuk dapat memenuhi kebutuhan peluang tersebut, diperlukan kebutuhan fungsional model sebagai berikut.

\begin{table}[h]
	\centering
	\caption{Spesifikasi Kebutuhan Fungsional Model}
	\label{tab:kebutuhan_fungsional}
	\begin{tabular}{|l|p{11cm}|}
		\hline
		\textbf{Kode} & \textbf{Kebutuhan Fungsional} \\ \hline
		FR01 & Model harus dapat menerima \textit{input} data berupa citra digital \textit{CT Scan} otak. \\ \hline
		FR02 & Model harus mampu mengklasifikasikan citra \textit{input} ke dalam dua kelas keputusan, yaitu tumor atau sehat (tidak ada tumor). \\ \hline
		FR03 & Model harus menghasilkan visualisasi interpretasi keputusan model untuk menyoroti area fitur pada citra yang membuat model mengklasifikasikan gambar citra tersebut ke dalam kelas tumor. \\ \hline
		FR04 & Model harus menampilkan tingkat kepercayaan dari hasil prediksi untuk membantu tenaga medis dalam menilai validitas keputusan sistem dan membangun kepercayaan pengguna. \\ \hline
	\end{tabular}
\end{table}

\subsection{Kebutuhan Nonfungsional}
Kebutuhan nonfungsional model adalah sebagai berikut.
\begin{table}[h]
	\centering
	\caption{Spesifikasi Kebutuhan Non-Fungsional Model}
	\label{tab:kebutuhan_non_fungsional}
	\begin{tabular}{|c|l|p{9.5cm}|}
		\hline
		\textbf{Kode} & \textbf{Kebutuhan} & \textbf{Deskripsi} \\ \hline
		NF01 & \textit{Interpretability} & Berbeda dengan model \textit{black-box} sebelumnya, model harus menyediakan transparansi atas keputusan yang diambil dengan memvisualisasikan area fokus sebagai dasar prediksi. Hal ini menjawab keraguan pasien terhadap validitas diagnosis AI dan membantu profesional medis memverifikasi hasil tanpa harus mempercayai algoritma secara buta. \\ \hline
		NF02 & \textit{Reliability} & Model harus memiliki konsistensi tinggi dalam memberikan hasil prediksi, dengan mempertahankan atau melampaui akurasi model referensi. \\ \hline
		NF03 & \textit{Complexity} & Model harus memiliki \textit{space complexity} yang rendah dan terukur pada tahap inferensi, agar dapat diimplementasikan secara optimal pada infrastruktur komputasi umum di fasilitas kesehatan tanpa mengorbankan kecepatan maupun akurasi klasifikasi dan visualisasi. \\ \hline
	\end{tabular}
\end{table}

\section{Analisis Pemilihan Solusi}
\label{sec:analisispemilihansolusi}
Pada bagian ini, dilakukan analisis kebutuhan untuk menentukan solusi yang paling tepat dalam mengatasi permasalahan \textit{black-box} pada model AI. Tahapan analisis ini mencakup identifikasi dan pemaparan berbagai alternatif solusi dan menggunakan metodologi AHP untuk mengevaluasi dan memilih solusi terbaik dari alternatif tersebut.

\subsection{Alternatif Solusi}

Masalah \textit{black-box} pada model AI sebenarnya berkaitan erat dengan konsep \textit{interpretability}, yaitu seberapa mudah manusia memahami alasan di balik suatu keputusan yang dihasilkan oleh model. \textit{Interpretability} dapat dicapai melalui dua cara. Cara pertama disebut \textit{post-hoc interpretability}, yaitu menambahkan metode penjelas pada model \textit{black-box} yang sudah ada tanpa mengubah model itu sendiri. Cara kedua disebut \textit{intrinsic interpretability}, yaitu menggunakan model yang sejak awal sudah mudah dipahami strukturnya, sehingga tidak memerlukan metode tambahan untuk menjelaskannya.

Berdasarkan dua cara tersebut, alternatif solusi pada tugas akhir ini juga dibagi menjadi dua. Alternatif pertama adalah menerapkan metode \textit{post-hoc interpretability} pada model \textit{deep learning}. Metode-metode ini sering disebut sebagai \textit{Explainable AI} (XAI). Metode XAI tidak benar-benar membuka isi model, melainkan hanya memberi perkiraan tentang bagian mana dari gambar yang paling berpengaruh terhadap keputusan model. Karena sifatnya perkiraan, hasil penjelasannya belum tentu 100\% sesuai dengan proses yang sebenarnya terjadi di dalam model. Penjelasan lebih lengkap mengenai metode-metode ini dapat dilihat pada Bab \ref{sec:metode_xai}, Subbab \nameref{sec:metode_xai}.

\begin{enumerate}
	\item \textit{Attribution-based methods} \\
	Metode ini bekerja dengan memberi nilai pada setiap bagian dari gambar masukan. Semakin besar pengaruh suatu bagian gambar terhadap keputusan akhir model, semakin tinggi nilai yang diberikan. Dengan cara ini, pengguna bisa melihat area mana pada gambar yang dianggap paling penting oleh model. Namun, karena metode ini hanya menghitung berdasarkan gradien setelah model selesai membuat keputusan, hasilnya tetap berupa perkiraan, bukan penjelasan langsung dari cara kerja model. Beberapa contoh metode \textit{attribution-based} yang dipertimbangkan adalah sebagai berikut.
	\begin{enumerate}
		\item \textit{Gradient-weighted Class Activation Mapping} (Grad-CAM)
		
		\item FullGrad-CAM dan FullGrad-CAM++
		
		\item \textit{SmoothGrad}
	\end{enumerate}
	
	\item \textit{Transformer-based methods} \\
	Metode ini memanfaatkan mekanisme \textit{attention} yang memang sudah ada di dalam arsitektur \textit{Vision Transformer} (ViT). Bedanya dengan \textit{attribution-based methods}, nilai \textit{attention} ini bukan hasil perkiraan tambahan, melainkan memang bagian dari proses perhitungan model itu sendiri. Oleh karena itu, metode ini bisa dibilang lebih dekat ke sifat \textit{intrinsic interpretability} dibanding metode berbasis gradien. Caranya adalah dengan membagi gambar menjadi beberapa \textit{patch}, lalu menampilkan skor \textit{attention}-nya dalam bentuk peta yang menunjukkan area mana yang paling penting bagi model \autocite{iftikhar2025explainable}.
\end{enumerate}

Selain mempertimbangkan berbagai metode \textit{post-hoc interpretability} di atas, tugas akhir ini juga mempertimbangkan alternatif kedua, yaitu mengganti model \textit{deep learning} yang digunakan dengan model \textit{machine learning} konvensional yang lebih mudah dipahami secara langsung, seperti \textit{decision tree} atau regresi logistik. Model semacam ini bersifat \textit{intrinsically interpretable} atau \textit{white-box}, karena proses pengambilan keputusannya bisa ditelusuri langsung dari strukturnya, tanpa perlu metode penjelas tambahan. Perbandingan antara pendekatan \textit{deep learning} yang dibantu \textit{post-hoc interpretability} dengan model konvensional (baik yang \textit{white-box} maupun yang masih \textit{black-box}) diperlukan untuk membuktikan solusi mana yang paling efektif dan cocok digunakan pada tugas akhir ini.

Untuk menentukan solusi mana yang akan dipilih, tugas akhir ini menggunakan metode \textit{Analytic Hierarchy Process} (AHP). AHP adalah metode pengambilan keputusan yang digunakan untuk memeringkat beberapa alternatif berdasarkan perbandingan berpasangan. Metode ini terdiri dari lima langkah utama. Pertama, menyusun masalah keputusan ke dalam tiga tingkatan, yaitu tujuan, kriteria, dan alternatif. Kedua, membandingkan setiap pasangan kriteria untuk menilai kepentingan relatifnya menggunakan skala sembilan poin, mulai dari 1 (sama penting) hingga 9 (sangat penting), beserta nilai kebalikannya. Ketiga, menghitung bobot masing-masing kriteria berdasarkan hasil perbandingan tersebut. Keempat, membandingkan setiap pasangan alternatif berdasarkan masing-masing kriteria. Kelima, mengalikan dan menjumlahkan bobot kriteria dengan skor alternatif untuk mendapatkan skor akhir setiap alternatif, yang kemudian digunakan untuk menentukan peringkatnya \autocite{1000minds2025ahp}.

AHP pada tugas akhir dilakukan berdasarkan literatur yang melakukan evaluasi terhadap model terkait. Terdapat tiga studi literatur yang digunakan untuk menganalisis alternatif solusi, yaitu:

\begin{enumerate}
	\item \textit{Explainable AI and vision transformers for detection and classification of brain tumor: a comprehensive survey} \cite{hosny2025explainable} \\
	Penelitian ini membahas tentang tinjauan teknis dan perbandingan teknik diagnostik \textit{deep learning} (DL) untuk tumor otak, dengan fokus pada penelitian dari tahun 2020 hingga 2024, seperti CNN, \textit{vision transformers} (ViT), teknik hibrida, dan \textit{Explainable AI} (XAI). Tugas akhir ini mengeksplorasi kinerja, keunggulan, dan keterbatasan dari pendekatan-pendekatan tersebut pada \textit{dataset} umum seperti Figshare dan BraTS.
	
	\item \textit{Explainable CNN for brain tumor detection and classification through XAI based key features identification} \cite{iftikhar2025explainable}\\
	Penelitian ini menggunakan metodologi yang menggabungkan \textit{Explainable AI} (XAI) dengan arsitektur CNN untuk mengatasi masalah model yang kompleks dan sulit diinterpretasi. Tugas akhir ini menggunakan teknik XAI berupa \textit{Gradient-weighted Class Activation Mapping} (Grad-CAM) agar memastikan model fokus pada fitur tumor yang relevan.
	
	\item \textit{Optimizing brain tumor classification through feature selection and hyperparameter tuning in machine learning models} \autocite{tahosin2023optimizing}\\
	Penelitian ini memberikan data perbandingan algoritma \textit{machine learning}. Hasilnya menunjukkan bahwa \textit{Extra Trees} mencapai akurasi tertinggi (98\%) dan \textit{Naive Bayes} memiliki akurasi terendah (69\%) namun sangat efisien secara komputasi.
\end{enumerate}

Tahapan AHP yang dilakukan adalah sebagai berikut:
\begin{enumerate}
	\item Struktur Hierarki Keputusan
	\begin{enumerate}
		\item Tujuan \\
		Memilih pendekatan \textit{deep learning} atau \textit{machine learning} yang paling efektif dan sesuai untuk tugas klasifikasi atau deteksi penyakit dari data medis.
		\item Kriteria
		\begin{enumerate}
			\item C1: Akurasi Deteksi \\
			Kemampuan model untuk secara benar mengidentifikasi atau mengklasifikasikan kondisi medis.
			\item C2: \textit{Interpretability} \\
			Kemampuan model untuk memberikan penjelasan atas prediksinya, yang krusial dalam konteks medis untuk membangun kepercayaan dan memahami fokus model.
			\item C3: Efisiensi Komputasi \\
			Kebutuhan sumber daya komputasi, waktu pelatihan, dan memori yang diperlukan model untuk mencapai performa yang diinginkan.
			\item C4: Kebutuhan Data dan Teknik \textit{Preprocessing}\\
			Sensitivitas model terhadap jumlah data yang tersedia dan efektivitas teknik seperti augmentasi data, \textit{transfer learning}, atau \textit{preprocessing} spesifik.
			\item C5: Kompleksitas Model \\
			Ukuran arsitektur, seperti jumlah \textit{layer}, parameter yang dapat memengaruhi interpretasi, efisiensi, dan risiko \textit{overfitting}.
		\end{enumerate}
		\item Alternatif Solusi
		\begin{enumerate}
			\item A1: CNN + Grad-CAM (Representasi \textit{attribution-based methods})
			\item A2: \textit{Vision Transformer} + \textit{Attention Map }(Representasi \textit{Transformer-based methods})
			\item A3: \textit{Naive Bayes} (Representasi \textit{white-box})
			\item A4: \textit{Extra Trees} (Representasi \textit{black-box})
		\end{enumerate}
	\end{enumerate}
	\item Perbandingan Berpasangan Antar Kriteria \\
	Perbandingan berpasangan dilakukan berdasarkan kriteria yang ada untuk menentukan prioritas dari setiap kriteria. Perbandingan antarkriteria dilakukan menggunakan matriks perbandingan berpasangan kriteria. Matriks perbandingan berpasangan kriteria adalah tabel yang merekam hasil perbandingan preferensi antara setiap pasang kriteria, menggunakan skala sembilan poin untuk menunjukkan tingkat kepentingan satu kriteria terhadap yang lain. Hasil dari matriks perbandingan kriteria dapat dilihat pada Tabel \ref{tab:perbandingan_kriteria}.
	\begin{table}[H]
		\centering
		\caption{Matriks Perbandingan Berpasangan Kriteria}
		\label{tab:perbandingan_kriteria}
		\begin{tabular}{|l|c|c|c|c|c|}
			\hline
			\textbf{Kriteria} & \textbf{C1} & \textbf{C2} & \textbf{C3} & \textbf{C4} & \textbf{C5} \\ \hline
			C1 (Akurasi) & $1$ & $1$ & $5$ & $3$ & $5$ \\ \hline
			C2 (XAI) & $1$ & $1$ & $5$ & $3$ & $5$ \\ \hline
			C3 (Efisiensi) & $1/5$ & $1/5$ & $1$ & $1/3$ & $1$ \\ \hline
			C4 (Data) & $1/3$ & $1/3$ & $3$ & $1$ & $3$ \\ \hline
			C5 (Kompleksitas) & $1/5$ & $1/5$ & $1$ & $1/3$ & $1$ \\ \hline
		\end{tabular}
	\end{table}
	Setelah dilakukan perbandingan, dilakukan normalisasi ke vektor prioritas kriteria. Normalisasi dilakukan dengan cara matriks perbandingan dikuadratkan, kemudian setiap baris dijumlahkan, dan hasil penjumlahan setiap baris dibagi dengan total keseluruhan jumlah baris. Langkah-langkah tersebut diulangi pada matriks hasil hingga nilai vektor prioritas yang dihasilkan menjadi konvergen \textcite{1000minds2025ahp}. Hasil dari vektor prioritas kriteria dapat dilihat pada Tabel \ref{tab:prioritas_kriteria}.
	\begin{table}[h]
		\centering
		\caption{Vektor Prioritas Kriteria}
		\label{tab:prioritas_kriteria}
		\begin{tabular}{|l|r|r|}
			\hline
			\textbf{Kriteria} & \textbf{Bobot} & \textbf{Prioritas} (\%) \\ \hline
			C1: Akurasi & 359 & 35,9\% \\ \hline
			C2: \textit{Interpretability} & 359 & 35,9\% \\ \hline
			C3: Efisiensi Komputasi & 64 & 6,4\% \\ \hline
			C4: Kebutuhan Data & 153 & 15,3\% \\ \hline
			C5: Kompleksitas Model & 64 & 6,4\% \\ \hline
			Total & 1.000 & 100,0\% \\ \hline
		\end{tabular}
	\end{table}
	\item Perbandingan Berpasangan Antar Alternatif \\
	Perbandingan dilakukan dengan cara setiap alternatif dibandingkan dengan alternatif solusi lain berdasarkan masing-masing kriteria menggunakan matriks perbandingan berpasangan kriteria \textcite{1000minds2025ahp}. Setelah itu, dilakukan normalisasi yang sama dengan perhitungan prioritas kriteria. Perbandingan antaralternatif berdasarkan kriteria dan prioritas alternatif solusi dapat dilihat pada Tabel \ref{tab:perbandingan_c1} hingga Tabel \ref{tab:perbandingan_c5}.
	\begin{table}[H]
		\centering
		\caption{Matriks Perbandingan Berpasangan C1 (Akurasi)}
		\label{tab:perbandingan_c1}
		\begin{tabular}{|l|c|c|c|c|}
			\hline
			\textbf{C1: Akurasi} & \textbf{A1} & \textbf{A2} & \textbf{A3} & \textbf{A4} \\ \hline
			\textbf{A1} & $1$ & $2$ & $7$ & $4$ \\ \hline
			\textbf{A2} & $1/2$ & $1$ & $5$ & $3$ \\ \hline
			\textbf{A3} & $1/7$ & $1/5$ & $1$ & $1/3$ \\ \hline
			\textbf{A4} & $1/4$ & $1/3$ & $3$ & $1$ \\ \hline
		\end{tabular}
	\end{table}
	
	Berdasarkan Tabel \ref{tab:perbandingan_c1}, dihasilkan pembobotan dengan menggunakan normalisasi sebagai berikut:
	\begin{enumerate}
		\item A1: 50,5\%
		\item A2: 30,6\%
		\item A3: 5,8\%
		\item A4: 13,1\%\\
	\end{enumerate}
	
	Pendekatan CNN dengan Grad-CAM (A1) memperoleh bobot tertinggi karena arsitektur CNN, terutama yang menggunakan \textit{transfer learning}, secara konsisten mendominasi performa akurasi dalam analisis citra medis dan terbukti lebih stabil pada \textit{dataset} standar. \textit{Vision Transformer} (A2) masih menghadapi tantangan generalisasi jika data latih terbatas, sedangkan \textit{Naive Bayes} (A3) memiliki performa paling rendah dengan akurasi sekitar 69\% yang dinilai kurang memadai untuk diagnosis kritis.
	
	\begin{table}[H]
		\centering
		\caption{Matriks Perbandingan Berpasangan C2 (\textit{Interpretability})}
		\label{tab:perbandingan_c2}
		\begin{tabular}{|l|c|c|c|c|}
			\hline
			\textbf{C2: XAI} & \textbf{A1} & \textbf{A2} & \textbf{A3} & \textbf{A4} \\ \hline
			\textbf{A1} & $1$ & $3$ & $2$ & $7$ \\ \hline
			\textbf{A2} & $1/3$ & $1$ & $1/3$ & $3$ \\ \hline
			\textbf{A3} & $1/2$ & $3$ & $1$ & $7$ \\ \hline
			\textbf{A4} & $1/7$ & $1/3$ & $1/7$ & $1$ \\ \hline
		\end{tabular}
	\end{table}
	
	Berdasarkan Tabel \ref{tab:perbandingan_c2}, dihasilkan pembobotan dengan menggunakan normalisasi sebagai berikut:
	\begin{enumerate}
		\item A1: 47,4\%
		\item A2: 13,9\%
		\item A3: 33,4\%
		\item A4: 5,3\%\\
	\end{enumerate}
	
	Metode CNN yang dilengkapi Grad-CAM (A1) unggul karena mampu menghasilkan visualisasi \textit{heatmap} yang menyoroti area tumor berdasarkan gradien lapisan konvolusi terakhir sehingga memberikan transparansi keputusan yang sangat dibutuhkan. \textit{Naive Bayes} (A3) menempati posisi kedua karena sifatnya sebagai model \textit{white-box} yang mudah dijelaskan secara matematis, berbeda dengan \textit{Extra Trees} (A4) yang merupakan model \textit{ensemble} (\textit{black-box}) yang sulit dilacak alur keputusannya. Sementara itu, peta atensi pada ViT (A2) sering kali menangkap hubungan global yang lebih abstrak dibandingkan fitur fisik tumor sehingga interpretasi visualnya lebih sulit dipahami dibandingkan \textit{heatmap} CNN.
	
	\begin{table}[H]
		\centering
		\caption{Matriks Perbandingan Berpasangan C3 (Efisiensi Komputasi)}
		\label{tab:perbandingan_c3}
		\begin{tabular}{|l|c|c|c|c|}
			\hline
			\textbf{C3: Efisiensi} & \textbf{A1} & \textbf{A2} & \textbf{A3} & \textbf{A4} \\ \hline
			\textbf{A1} & $1$ & $3$ & $1/7$ & $1/3$ \\ \hline
			\textbf{A2} & $1/3$ & $1$ & $1/9$ & $1/5$ \\ \hline
			\textbf{A3} & $7$ & $9$ & $1$ & $3$ \\ \hline
			\textbf{A4} & $3$ & $5$ & $1/3$ & $1$ \\ \hline
		\end{tabular}
	\end{table}
	
	Berdasarkan Tabel \ref{tab:perbandingan_c3}, dihasilkan pembobotan dengan menggunakan normalisasi sebagai berikut:
	\begin{enumerate}
		\item A1: 10,1\%
		\item A2: 4,9\%
		\item A3: 60,7\%
		\item A4: 24,3\%\\
	\end{enumerate}
	
	\textit{Naive Bayes} (A3) mendominasi kriteria ini karena kesederhanaan algoritmanya yang membutuhkan sumber daya komputasi minimal dan waktu pelatihan yang cepat. \textit{Extra Trees} (A4) juga dinilai efisien karena struktur \textit{tree}-nya yang mengurangi redundansi data. Metode berbasis \textit{deep learning} seperti CNN (A1) dan \textit{Vision Transformer} (A2) memiliki beban komputasi yang jauh lebih berat, seperti ViT yang memproses gambar sebagai urutan \textit{patch} dengan mekanisme atensi yang memakan memori besar dan waktu pelatihan yang jauh lebih lama dibandingkan metode konvensional.
	
	\begin{table}[H]
		\centering
		\caption{Matriks Perbandingan Berpasangan C4 (Kebutuhan Data)}
		\label{tab:perbandingan_c4}
		\begin{tabular}{|l|c|c|c|c|}
			\hline
			\textbf{C4: Data} & \textbf{A1} & \textbf{A2} & \textbf{A3} & \textbf{A4} \\ \hline
			\textbf{A1} & $1$ & $3$ & $1/5$ & $1/3$ \\ \hline
			\textbf{A2} & $1/3$ & $1$ & $1/7$ & $1/5$ \\ \hline
			\textbf{A3} & $5$ & $7$ & $1$ & $2$ \\ \hline
			\textbf{A4} & $3$ & $5$ & $1/2$ & $1$ \\ \hline
		\end{tabular}
	\end{table}
	
	Berdasarkan Tabel \ref{tab:perbandingan_c4}, dihasilkan pembobotan dengan menggunakan normalisasi sebagai berikut:
	\begin{enumerate}
		\item A1: 12,2\%
		\item A2: 5,7\%
		\item A3: 52,3\%
		\item A4: 29,8\%\\
	\end{enumerate}
	
	Alternatif \textit{Naive Bayes} (A3) dan \textit{Extra Trees} (A4) lebih unggul karena kemampuan mereka untuk bekerja efektif dan tidak mudah \textit{overfitting} pada \textit{dataset} dengan jumlah sampel yang sedikit. Sebaliknya, CNN (A1) membutuhkan jumlah data yang besar atau strategi augmentasi data untuk mencapai generalisasi yang baik, dan \textit{Vision Transformer} (A2) mendapatkan bobot terendah karena arsitekturnya tidak memiliki bias induktif sehingga memerlukan \textit{dataset} berskala besar agar performanya tidak menurun drastis.
	
	\begin{table}[H]
		\centering
		\caption{Matriks Perbandingan Berpasangan C5 (Kompleksitas)}
		\label{tab:perbandingan_c5}
		\begin{tabular}{|l|c|c|c|c|}
			\hline
			\textbf{C5: Kompleksitas} & \textbf{A1} & \textbf{A2} & \textbf{A3} & \textbf{A4} \\ \hline
			\textbf{A1} & $1$ & $2$ & $1/7$ & $1/3$ \\ \hline
			\textbf{A2} & $1/2$ & $1$ & $1/8$ & $1/4$ \\ \hline
			\textbf{A3} & $7$ & $8$ & $1$ & $3$ \\ \hline
			\textbf{A4} & $3$ & $4$ & $1/3$ & $1$ \\ \hline
		\end{tabular}
	\end{table}
	
	Berdasarkan Tabel \ref{tab:perbandingan_c5}, dihasilkan pembobotan dengan menggunakan normalisasi sebagai berikut:
	\begin{enumerate}
		\item A1: 9,4\%
		\item A2: 6,0\%
		\item A3: 60,8\%
		\item A4: 23,8\%\\	
	\end{enumerate}
	
	Dalam aspek kompleksitas arsitektur, \textit{Naive Bayes} (A3) dinilai paling sederhana karena hanya mengandalkan perhitungan probabilitas statistik dasar, diikuti oleh \textit{Extra Trees} (A4). Sementara itu, CNN (A1) memiliki struktur hierarkis yang dalam dengan jutaan parameter yang harus dilatih, dan \textit{Vision Transformer} (A2) dinilai paling kompleks karena mekanisme \textit{multi-head self-attention} yang memiliki kompleksitas kuadratik terhadap jumlah token (\textit{patch} gambar), menjadikannya model yang paling berat secara struktural dibandingkan alternatif lainnya.
	
	\item Perbandingan Berpasangan Setiap Pasangan Alternatif Berdasarkan Setiap Kriteria \\
	Setelah perbandingan dilakukan pada setiap alternatif berdasarkan kriteria, dilakukan penggabungan perbandingan pada setiap pasangan alternatif berdasarkan setiap kriteria dengan hasil yang dapat dilihat pada Tabel \ref{tab:hasil_perhitungan}.
	\begin{table}[h]
		\centering
		\caption{Hasil Perhitungan Perbandingan Berpasangan}
		\label{tab:hasil_perhitungan}
		\begin{tabular}{|l|c|c|c|c|c|c|}
			\hline
			\textbf{Alternatif} & \textbf{C1} & \textbf{C2} & \textbf{C3} & \textbf{C4} & \textbf{C5} & \textbf{Hasil Akhir} \\ \hline
			\textbf{Bobot Kriteria} & 35,9\% & 35,9\% & 6,4\% & 15,3\% & 6,4\% & \\ \hline
			\textbf{A1} & 505 & 474 & 101 & 122 & 94 & 382,607 \\ \hline
			\textbf{A2} & 306 & 139 & 49 & 57 & 60 & 175,452 \\ \hline
			\textbf{A3} & 58 & 334 & 607 & 523 & 608 & 298,507 \\ \hline
			\textbf{A4} & 131 & 53 & 243 & 298 & 238 & 142,434 \\ \hline
		\end{tabular}
	\end{table}
	
	Perhitungan Rinci Skor Akhir:
	\begin{enumerate}
		\item A1 (CNN + Grad-CAM): $(0,359 \times 505) + (0,359 \times 474) + (0,064 \times 101) + (0,153 \times 122) + (0,064 \times 94) = 382,607$
		\item A2 (ViT + \textit{Attention Map}): $(0,359 \times 306) + (0,359 \times 139) + (0,064 \times 49) + (0,153 \times 57) + (0,064 \times 60) = 175,452$
		\item A3 (\textit{Naive Bayes}): $(0,359 \times 58) + (0,359 \times 334) + (0,064 \times 607) + (0,153 \times 523) + (0,064 \times 608) = 298,507$
		\item A4 (\textit{Extra Trees}): $(0,359 \times 131) + (0,359 \times 53) + (0,064 \times 243) + (0,153 \times 298) + (0,064 \times 238) = 142,434$
	\end{enumerate}
	\item Pemilihan Solusi \\
	Dengan mempertimbangkan bobot prioritas dan skor kinerja keseluruhan, model \textit{deep learning} yang menggunakan Grad-CAM (A1) ditetapkan sebagai solusi pada tugas akhir. Metode ini memberikan keseimbangan terbaik antara akurasi tinggi dan interpretabilitas kuat, sejalan dengan tujuan metodologi yang menekankan pada keandalan model sekaligus transparansi keputusan AI.
\end{enumerate}